% Simulation Settings (INFOCOM-style run-in paragraphs).
% Numbers match sim/config.py + run_v2x_real.py overrides and the
% Table-I runs (metrics_v2x_real_{kitti,nuscenes}.npz, 3 seeds).
\subsection{Simulation Settings}
\label{subsec:settings}

\textbf{Mobility trace and V2V network.}
We evaluate FACE on a real vehicular mobility trace of the
Seoul--Gangnam district, the operating area of the robotaxi service in
Sec.~II, reconstructed over the real road network (38{,}091 road
segments) from open Seoul traffic data. The trace contains
$|\mathcal I|=180$ vehicles sampled at $\Delta_{\mathrm{rd}}\approx10$\,s;
training runs for $T=250$ rounds by replaying the trace window
cyclically (steady-state traffic). Vehicles communicate within
$r_{\mathrm{V2V}}=150$\,m, a transmission succeeds with probability
$1-(d_{ij}/r_{\mathrm{V2V}})^2$, and each contact provides a budget
$\widehat B_{ij}$ of $1.6$\,s at an effective sidelink rate of
$12$\,MB/s, so a typical contact carries roughly one encoder. Road
zones $\mathcal Z$ are $300$\,m grid cells over the district.

\textbf{Datasets and models.}
Following the multimodal perception task of Sec.~II, each vehicle
trains a multimodal object classifier on real sensor data: KITTI
(camera $+$ LiDAR crops of labeled objects; Car / Pedestrian /
Cyclist) and nuScenes (camera $+$ LiDAR $+$ radar; Car / Pedestrian). Modality availability is
heterogeneous: $\Pr(\text{LiDAR})=0.85$ and, on nuScenes,
$\Pr(\text{radar})=0.7$, so a share of vehicles is camera-only. Each
vehicle holds $80$--$600$ local training objects with per-vehicle
sensing-quality degradation; $15\%$ of the vehicles carry clean
high-quality data and act as encoder carriers, while the remaining
vehicles form the high-demand cohort reported in Table~I. Encoders
are modality-specific networks with a local fusion head; encoder
transmission sizes are $S_x=12/18/6$\,MB (camera/LiDAR/radar),
representative of production perception encoders. Aggregation follows
Sec.~III with data-size weights as in FedAvg~\cite{fedavg}.

\textbf{FACE parameters.}
Unless stated otherwise, the future-contact horizon is $H=6$ rounds
with discount $\beta=0.85$, the communication price is
$\lambda=10^{-3}$ per MB, the copy cap is $K_x=6$ with version
lifetime $15$ rounds and republication every round, the adoption
threshold is $\delta=5\times10^{-3}$, the evaluation budget is
$N_{\mathrm{ev}}=6$, and the cache capacity is
$\Lambda_i^{\max}=45$\,MB. The gain predictor uses a sliding-window
ridge (regularization $1.0$, decay $0.98$); zone statistics decay with
factor $0.97$ and forecast blend $\gamma_f=0.85$. All results average
three independent runs (distinct seeds for data assignment, modality
availability, and matching); experiments run on a single NVIDIA RTX
3080 GPU.

\textbf{Benchmarks.}
All schemes run under the identical Sec.~III protocol (encoder types
and versions, replication tokens, evaluation-gated adoption) and
differ only in their forwarding and caching policy, so the comparison
isolates the delivery decisions:
(1)~\textbf{Cached-DFL}~\cite{cachedDFL}: delay-tolerant
model caching on mobile agents -- each vehicle caches encoders
received from encountered vehicles and re-shares its own and cached
encoders at every contact with LRU replacement, without demand
awareness;
(2)~\textbf{V2V}: each vehicle opportunistically shares its own
encoders with the neighbors offering the highest instantaneous V2V
link quality, without considering encoder demand;
(3)~\textbf{Learning}: each vehicle shares its own encoders directly
with the neighbors that currently demand them, prioritizing receivers
by expected learning benefit regardless of link quality, and performs
no encoder caching for subsequent relaying;
(4)~\textbf{mmFedMC}~\cite{multimodal_tmc}: at each contact, the
sender offers only its highest-value own-modality encoder, ranked by
learning contribution per unit communication cost, following the
modality selection of mmFedMC adapted to the decentralized setting;
(5)~\textbf{AutoFed}~\cite{autofed}: each vehicle offers encoders to
compatible receivers in descending order of encoder quality, without
considering receiver demand.

\textbf{Metrics.}
We report the final test accuracy (mean of the last 20 rounds
$\pm$ std over runs) for all vehicles and for the high-demand cohort,
the validation loss, the useful-delivery ratio of Sec.~III, the
deadline-delivery probability Delivery@$d$ (a useful encoder reaches a
high-demand vehicle within $d$ rounds), and the rounds and the
communication volume Comm@$\tau$ (cumulative encoder bytes in GB,
matching the cost $C(t)$ of Sec.~III) to reach the target accuracy
$\tau$, set to $95\%$ of the best final accuracy on each dataset.
